\documentclass[10pt,a4paper]{amsart}
\usepackage[margin=19mm]{geometry}
\usepackage{amsmath,amssymb,mathtools}
\usepackage[scr=rsfs,cal=euler]{mathalfa}
\usepackage{xfrac}
\usepackage[unicode,hidelinks,bookmarks=false]{hyperref}
\newcommand{\ie}{i.e.\ }
\newcommand{\cf}{cf.\ }
\newcommand{\resp}{respectively\ }
\newcommand{\Subsection}{\subsection}
\providecommand{\nobreakdash}{\nobreak\hbox{-}}
\setlength{\emergencystretch}{2em}
\multlinegap0pt
\usepackage{bm}
\usepackage{amssymb}
\usepackage{enumerate}
\usepackage{algorithm}\usepackage{algpseudocode}
\newtheorem{lemma}{Lemma}
\newtheorem{theorem}{Theorem}
\newcommand{\LRp}[1]{\left( #1 \right)}
\newcommand{\algn}[1]{\begin{align} #1 \end{align}}
\newcommand{\algns}[1]{\begin{align*} #1 \end{align*}}
\newcommand{\ta}{\tau}
\title{Optimal control of poroelasticity equations}
\author{Arbaz Khan, Jeonghun J. Lee, Harpal Singh}
\date{Source: arXiv:2605.30839v1 (CC BY 4.0)}
\begin{document}
\maketitle

\noindent\emph{Source and licence.} Optimal control of poroelasticity equations. Version arXiv:2605.30839v1 (CC BY 4.0), distributed by arXiv under the Creative Commons Attribution 4.0 International licence, http://creativecommons.org/licenses/by/4.0/, as recorded in the arXiv licence metadata for this exact version and shown on the abstract page https://arxiv.org/abs/2605.30839v1. Full licence notice: \texttt{licenses/arxiv-2605.30839-CC-BY-4.0.txt}.\par\smallskip

\noindent\textit{Editorial scope.} Counted unit: the article's theorem Theorem:5.9 (source0.tex lines 2417--2461). The article's complete original proof of it is delivered with it (source0.tex lines 2463--2604, one \texttt{proof} environment of 142 lines). No mathematics is written, repaired or re-derived: the statement and the proof are the article's own lines, in the article's own notation, and no retained line is substituted or re-flowed.  Retained uncounted as the source-local context the proof needs: lines 752--759; lines 1052--1084; lines 1168--1171; lines 1209--1238; lines 1962--1993; lines 2321--2337.  Nothing was omitted from any retained span, so the package carries no omission record.  No retained line contains a citation, so the wrapper carries no bibliography and no bibliography entry.  No retained line contains a figure environment, an image-inclusion command or drawing code, so no image is delivered and the package declares no asset dependency.  Editorial formatting only, in addition: an AMS \texttt{amsart} preamble that restates the article's own declarations and macros used by the retained lines, namely the environments \texttt{lemma}, \texttt{theorem} on separate counters, together with the standard packages those lines need.  The retained environments print in this document as Lemma 1, Theorem 1 in order of appearance, whereas the article prints the same items with the numbering of its own full document.  The complete native source of the article as distributed in the arXiv e-print for this exact version is bundled unchanged as \texttt{sources/201695/source0.tex}.
\begin{align}
	\label{obj_fun} J \LRp{\boldsymbol{u}, p, m} = %\frac{\alpha_{\boldsymbol{u}, T}}{2} \int_{\Omega} |\boldsymbol{u} \LRp{x, T} -  \boldsymbol{u}_T (x)|^2 \, dx 
   % \\ \notag
	  \frac{\alpha_{\boldsymbol{u}, C}}{2} \int_0^T \int_{\Omega} |\boldsymbol{u} \LRp{x, t} \hspace{-0.2mm}-\hspace{-0.2mm} \boldsymbol{u}_C \LRp{x, t}|^2 \, dx \, dt 
     %+ \frac{\alpha_{p, T}}{2} \int_{\Omega} |p\LRp{x, T} - p_T (x)|^2 \, dx 
     \\
	 + \frac{\alpha_{p, C}}{2} \int_0^T \int_{\Omega} |p(x, t) - p_C (x, t)|^2 \, dx \, dt + \frac{\gamma}{2} \int_0^T \int_{\Omega} |m(x, t)|^2 \, dx \, dt \notag,
\end{align}

\medskip
{\bf Discretization of forward equation} The numerical solution at the $(k+1)-$th time step is defined inductively as follows: For $0\le k \le N-1$ find $(U^{k+1}, P^{k+1}) \in \boldsymbol{V}_h \times Q_h$ such that
%
\begin{subequations}\label{eq:fullydiscrete-eqs}
	\algn{
		\label{eq:fullydiscrete-eq1}
		a_{\boldsymbol{u}} \LRp{\partial_\ta U^{k+1}, \boldsymbol{v}} + b \LRp{\partial_\ta P^{k+1}, \boldsymbol{v}} &= \frac{1}{\Delta t} \LRp{{\boldsymbol{f}}, \boldsymbol{v}}_{I_k\times\Omega}, \\
		\label{eq:fullydiscrete-eq2}
		b \LRp{q, \partial_\ta U^{k+1}} - \LRp{s_0 \partial_\ta P^{k+1}, q}_{\Omega} - a_p\LRp{ P^{k+1}, q} &= \frac{1}{\Delta t}\LRp{m, q}_{I_k\times\Omega}, 
	}    
\end{subequations}
for all $\LRp{\boldsymbol{v}, q} \in \boldsymbol{V}_h \times Q_h$. 
%
\begin{lemma}
    The system \eqref{eq:fullydiscrete-eqs} has a unique solution.     
\end{lemma}
%
\begin{proof}
To establish the well-posedness of \eqref{eq:fullydiscrete-eqs}, we show that $(U^{k+1}, P^{k+1})$ is uniquely determined by the linear system (\ref{eq:fullydiscrete-eqs})  when $U^{k}, P^{k}, {\boldsymbol{f}}$ and $m$ are given. Equivalently, it is enough to show that $U^{k+1}=0$,  $P^{k+1}=0$ when $U^{k}, P^{k}, \boldsymbol{f}, m$ are zero. If we assume that $U^{k}, P^{k}, \boldsymbol{f}, m$ are zero, \eqref{eq:fullydiscrete-eqs} can be written as:
\begin{subequations}
	\algns{
		a_{\boldsymbol{u}} \LRp{U^{k+1}, \boldsymbol{v}} + b \LRp{P^{k+1}, \boldsymbol{v}} = 0,
		\\
		b \LRp{q, U^{k+1}} - \LRp{s_0 P^{k+1}, q}_{\Omega} - \Delta t\LRp{{\kappa} \nabla P^{k+1}, \nabla q}_{\Omega} = 0, 
	}    
\end{subequations}
%
for all $\LRp{\boldsymbol{v}, q} \in \boldsymbol{V}_h \times Q_h$. Choosing $ \boldsymbol{v}=U^{k+1}$ and $q=-P^{k+1}$ and adding both equations, we have
\begin{align*}
	0\le  a_{\boldsymbol{u}} \LRp{U^{k+1}, U^{k+1}}+  \LRp{s_0 P^{k+1}, P^{k+1}}_{\Omega} +\Delta t a_p\LRp{ P^{k+1}, P^{k+1}}=0.
\end{align*}
This implies that $U^{k+1}=0 = P^{k+1}$.    
\end{proof}

\begin{align}
    \label{eq:initial-data-assumption}
    \|\boldsymbol{u}(0)-U^0\|_{H^1(\Omega)}+\|p(0)-P^0\|_{H^1(\Omega)}\le C h^{s}(\|\boldsymbol{u}(0)\|_{H^{1+s}(\Omega)}+\|p(0)\|_{H^{1+s}(\Omega)}), \quad s >0.
\end{align}

\begin{lemma}\label{lemma:aux-forward-error}
    %Assume that {\rm (D)} and {\rm (E)} hold with $s_{\min} = \min \{s_e, s_d\}>0$. 
	Let $({\boldsymbol{u}},{p},{m})$ be the solution of the optimal control problem \eqref{obj_fun}. Suppose that $(U({m}), P({m})) \in \boldsymbol{V}_{hk} \times Q_{hk}$ is a solution of \eqref{eq:fullydiscrete-eqs} for the continuous optimal control ${m}$ with numerical initial data $(U^0, P^0)$ satisfying \eqref{eq:initial-data-assumption} for some $0<s\le 1$. Assume also that 
    %
    \[
        \boldsymbol{u}\in H^1(0,T;\boldsymbol{H}^{1+s}(\Omega)), \quad  p\in H^1(0,T;H^{1+s}(\Omega)).
    \]
    %
    Then, 
    %
	\begin{align}
        \label{eq:forward-aux-estimate}
		&\max_{1\le k \le N} \LRp{\| \boldsymbol{u}^k - U^k(m)\|_{H^1(\Omega)} + \| p^k - P^k(m)\|_{L^2(\Omega)} } 
        \\
        \notag
		&+ \| \boldsymbol{u} - U(m)\|_{L^2(0,T; H^1(\Omega))} + \| p - P(m)\|_{L^2(0,T; L^2(\Omega))} 
        \\
        \notag
        &\le h^{s}(\|\boldsymbol{u}(0)\|_{H^{1+s}(\Omega)}+\|p(0)\|_{H^{1+s}(\Omega)}) 
        \\
        \notag
        &\quad + C(h^s+\Delta t)
        \left(
        \|\boldsymbol{u}\|_{H^1(0,T;H^{1+s}(\Omega))}
        +
        \|p\|_{H^1(0,T;H^{1+s}(\Omega))}
        \right) .
    \end{align}
    %
\end{lemma}

\begin{lemma}[adjoint error for a fixed control]
\label{lem:adjoint-error-fixed-control}
Let $m\in \mathcal{M}_{ad}$ be fixed. Let $(\boldsymbol{u},p)$ be the
corresponding continuous state and let $(\boldsymbol{w},r)$ be the
corresponding continuous adjoint. Let $(U(m),P(m))$ be the fully
discrete state generated by the same control $m$, and let
$(W(m),R(m))$ be the corresponding fully discrete adjoint.

Assume that
\[
 r\in H^1(0,T;H^{1+s}(\Omega)),
 \quad 
 \boldsymbol{w}\in H^1(0,T;\boldsymbol{H}^{1+s}(\Omega)),
\]
for some $0<s\le 1$. Then
%
\begin{multline}
    \|R(m)-r\|_{L^2(0,T;L^2(\Omega))} 
    \\
    \le C\left( h^s+\Delta t + \|U(m)-\boldsymbol{u}\|_{L^2(0,T;L^2(\Omega))} +\|P(m)-p\|_{L^2(0,T;L^2(\Omega))} \right).
\end{multline}
%
Consequently, if the fixed-control state error satisfies
\[
\|U(m)-\boldsymbol{u}\|_{L^2(0,T;L^2(\Omega))}
+ \|P(m)-p\|_{L^2(0,T;L^2(\Omega))} \le C(h^s+\Delta t),
\]
then
\[
\|R(m)-r\|_{L^2(0,T;L^2(\Omega))} \le C(h^s+\Delta t).
\]
\end{lemma}

\begin{theorem}\label{Theorem:5.8}
	Let $(\boldsymbol{u},p,m)$ and $(U,P,M)$ be the admissible solutions of the continuous and fully discrete optimal control problems. Suppose that $(\boldsymbol{u}, p)$ and the continuous adjoint solution satisfy
	\begin{align}\label{eq:adjoint-regularity}
		p, r \in H^1(0,T; H^{1+s}(\Omega)) ,
		\qquad
		\boldsymbol{u}, \boldsymbol{w} \in H^1(0,T; \boldsymbol{H}^{1+s}(\Omega)),
	\end{align}
    %
    and the initial data assumption \eqref{eq:initial-data-assumption} holds. 
	Then, we have the estimate 
	%
    \begin{align*}
		\|m-M\|_{L^2(0,T;L^2(\Omega))} &\le C (\Delta t + h^s)
	\end{align*}
    %
    with $C>0$ depending on the solution regularities in \eqref{eq:adjoint-regularity}.
\end{theorem}

\begin{theorem}\label{Theorem:5.9}
Let $(\boldsymbol u,p,\bar m)$ be the continuous optimal solution and
let $(U,P,M)$ be the fully discrete optimal solution. Under the
regularity assumptions of Lemma~\ref{lemma:aux-forward-error},
Lemma~\ref{lem:adjoint-error-fixed-control}, and
Theorem~\ref{Theorem:5.8}, there exists a constant $C>0$, independent of
$h$ and $\Delta t$, such that
\[
\begin{aligned}
&\max_{1\le k\le N}
\left(
\|\boldsymbol u^k-U^k\|_{H^1(\Omega)}
+
\|p^k-P^k\|_{L^2(\Omega)}
\right)
\\
&\qquad
+
\|\boldsymbol u-U\|_{L^2(0,T;H^1(\Omega))}
+
\|p-P\|_{L^2(0,T;L^2(\Omega))}
\le
C(h^s+\Delta t).
\end{aligned}
\]
In particular, if $s=1$, then
\[
\begin{aligned}
&\max_{1\le k\le N}
\left(
\|\boldsymbol u^k-U^k\|_{H^1(\Omega)}
+
\|p^k-P^k\|_{L^2(\Omega)}
\right)
\\
&\qquad
+
\|\boldsymbol u-U\|_{L^2(0,T;H^1(\Omega))}
+
\|p-P\|_{L^2(0,T;L^2(\Omega))}
\le
C(h+\Delta t).
\end{aligned}
\]
\end{theorem}

\begin{proof}
Let $(U(\bar m),P(\bar m))$ denote the fully discrete state generated by
the continuous optimal control $\bar m$. We decompose the total state
errors as
\begin{align*}
\boldsymbol u^k-U^k
&=
\bigl(\boldsymbol u^k-U^k(\bar m)\bigr)
+
\bigl(U^k(\bar m)-U^k\bigr),
\\
p^k-P^k
&=
\bigl(p^k-P^k(\bar m)\bigr)
+
\bigl(P^k(\bar m)-P^k\bigr).
\end{align*}
The first terms are fixed-control discretization errors. By
Lemma~\ref{lemma:aux-forward-error},
\[
\begin{aligned}
&\max_{1\le k\le N}
\left(
\|\boldsymbol u^k-U^k(\bar m)\|_{H^1(\Omega)}
+
\|p^k-P^k(\bar m)\|_{L^2(\Omega)}
\right)
\\
&\qquad
+
\|\boldsymbol u-U(\bar m)\|_{L^2(0,T;H^1(\Omega))}
+
\|p-P(\bar m)\|_{L^2(0,T;L^2(\Omega))}
\le
C(h^s+\Delta t).
\end{aligned}
\]

It remains to estimate the perturbation caused by replacing $\bar m$ by
the discrete optimal control $M$. Set
\[
\zeta_{\boldsymbol u}^k:=U^k(\bar m)-U^k(M),
\qquad
\zeta_p^k:=P^k(\bar m)-P^k(M).
\]
Subtracting the two fully discrete state equations corresponding to
$\bar m$ and $M$, we obtain, for all
$(\boldsymbol v,q)\in \boldsymbol V_h\times Q_h$,
\begin{align}
a_{\boldsymbol u}(\partial_\tau\zeta_{\boldsymbol u}^k,\boldsymbol v)
+
b(\partial_\tau\zeta_p^k,\boldsymbol v)
&=0,
\\
b(q,\partial_\tau\zeta_{\boldsymbol u}^k)
-
(s_0\partial_\tau\zeta_p^k,q)_\Omega
-
a_p(\zeta_p^k,q)
&=
\frac1{\Delta t}
(\bar m-M,q)_{I_k\times\Omega}.
\end{align}
The initial data are the same for both discrete states, hence
$\zeta_{\boldsymbol u}^0=0$, $\zeta_p^0=0$.
For the right-hand side $G_m^k(q):=\frac1{\Delta t}(\bar m-M,q)_{I_k\times\Omega}$,
we have
\[
\|G_m^k\|_{a_p,*}
\le
C\Delta t^{-1/2}
\|\bar m-M\|_{L^2(I_k;L^2(\Omega))}.
\]
Consequently,
\[
\Delta t\sum_{k=0}^{N-1}\|G_m^k\|_{a_p,*}^2
\le
C\|\bar m-M\|_{L^2(0,T;L^2(\Omega))}^2.
\]
Applying the same stability argument as in the proof of
Lemma~\ref{lemma:aux-forward-error}, with zero consistency error and
with right-hand side $\bar m-M$, gives
\[
\begin{aligned}
&\max_{1\le k\le N}
\left(
\|\zeta_{\boldsymbol u}^k\|_{H^1(\Omega)}
+
\|\zeta_p^k\|_{L^2(\Omega)}
\right)
+
\|\zeta_{\boldsymbol u}\|_{L^2(0,T;H^1(\Omega))}
+
\|\zeta_p\|_{L^2(0,T;L^2(\Omega))}
\\
&\qquad
\le
C
\|\bar m-M\|_{L^2(0,T;L^2(\Omega))}.
\end{aligned}
\]
By Theorem~\ref{Theorem:5.8}, $\|\bar m-M\|_{L^2(0,T;L^2(\Omega))} \le C(h^s+\Delta t)$. 
Therefore,
\[
\begin{aligned}
&\max_{1\le k\le N}
\left(
\|\zeta_{\boldsymbol u}^k\|_{H^1(\Omega)}
+
\|\zeta_p^k\|_{L^2(\Omega)}
\right)
+
\|\zeta_{\boldsymbol u}\|_{L^2(0,T;H^1(\Omega))}
+
\|\zeta_p\|_{L^2(0,T;L^2(\Omega))}
\le
C(h^s+\Delta t).
\end{aligned}
\]

Combining the fixed-control discretization estimate and the control
perturbation estimate by the triangle inequality yields
\[
\begin{aligned}
&\max_{1\le k\le N}
\left(
\|\boldsymbol u^k-U^k\|_{H^1(\Omega)}
+
\|p^k-P^k\|_{L^2(\Omega)}
\right)
\\
&\qquad
+
\|\boldsymbol u-U\|_{L^2(0,T;H^1(\Omega))}
+
\|p-P\|_{L^2(0,T;L^2(\Omega))}
\le
C(h^s+\Delta t).
\end{aligned}
\]
For $s=1$, this gives the asserted $O(h+\Delta t)$ estimate.
\end{proof}

\end{document}
